\documentclass[10pt, aspectratio=169]{beamer}

% Configurações de Tema e Cores
\usetheme{Madrid}
\usecolortheme{beaver} % Esquema de cores vermelho/cinza (sóbrio para saúde)
\usefonttheme{professionalfonts}

% Pacotes Adicionais
\usepackage[utf8]{inputenc}
\usepackage[portuguese]{babel}
\usepackage{booktabs} % Para tabelas bonitas
\usepackage{graphicx}
\usepackage{natbib}
\usepackage{hyperref}
\usepackage{ragged2e} % Para justificar texto

% Configuração de Bibliografia no Beamer
\setbeamertemplate{bibliography item}{\insertbiblabel}

% Metadados
\title{Aprendizado de Máquina na Vigilância de Precisão da Hanseníase no Brasil}
\subtitle{Modelagem de Subnotificação Pandêmica, Reabsorção de Backlog e Explicabilidade SHAP}
\author{Aureziano Faria de Oliveira}
\institute[UFES]{Programa de Pós-Graduação em Informática \\ Universidade Federal do Espírito Santo (UFES)}
\date{Defesa de Dissertação --- 2026}

\begin{document}

% Slide de Título
\begin{frame}
    \titlepage
\end{frame}

% Sumário
\begin{frame}{Roteiro da Apresentação}
    \tableofcontents
\end{frame}

% ============================================================================
% CAPÍTULO 1: INTRODUÇÃO
% ============================================================================
\section{Introdução e Motivação}

\begin{frame}{O Cenário da Hanseníase no Brasil}
    \begin{columns}
        \column{0.6\textwidth}
        \begin{itemize}
            \item \textbf{Carga da Doença:} Brasil como 2º país com maior número de casos novos no mundo.
            \item \textbf{O Problema:} Diagnóstico tardio correlacionado ao Grau 2 de Incapacidade Física (G2D).
            \item \textbf{Impacto da Pandemia:} Interrupção da busca ativa e vigilância, gerando uma "Pandemia Oculta" \citep{who_hans_2021}.
            \item \textbf{Backlog:} Acúmulo de casos não detectados (2020-2022) que agora pressionam o sistema de saúde \citep{dahbeer2022impact}.
        \end{itemize}
        
        \column{0.4\textwidth}
        \begin{alertblock}{Risco de Subnotificação}
            A queda aparente na incidência mascara o aumento da gravidade real e a demanda reprimida.
        \end{alertblock}
    \end{columns}
\end{frame}

\begin{frame}{Objetivos da Pesquisa}
    \textbf{Objetivo Geral:} Analisar o impacto da pandemia e desenvolver modelos preditivos de risco.
    \vspace{0.5cm}
    \begin{enumerate}
        \item \textbf{Quantificar o "Apagão":} Usar séries temporais para medir o gap de notificação (2020-2024).
        \item \textbf{Avaliar Qualidade de Dados:} Medir o impacto de dados ausentes no SINAN (ex: Baciloscopia).
        \item \textbf{Modelagem Preditiva:} Comparar Random Forest, XGBoost e LightGBM na predição de G2D.
        \item \textbf{Mitigação de Viés:} Aplicar técnicas de balanceamento (SMOTE) para detectar a classe minoritária (casos graves).
    \end{enumerate}
\end{frame}

% ============================================================================
% CAPÍTULO 2: DADOS E EDA
% ============================================================================
\section{Preparação dos Dados e Análise Exploratória}

\begin{frame}{Caracterização do Dataset (SINAN)}
    \begin{itemize}
        \item \textbf{Volume de Dados:} $N = 998.214$ notificações (2001--2024).
        \item \textbf{Pré-processamento:}
            \begin{itemize}
                \item Remoção de colunas constantes (\texttt{TP\_NOT}, \texttt{ID\_AGRAVO}).
                \item Discretização da idade (Foco em $<15$ anos).
            \end{itemize}
        \item \textbf{Desafio dos Dados Ausentes (Missing Values):}
    \end{itemize}
    
    \begin{table}
        \centering
        \begin{tabular}{lcc}
            \toprule
            \textbf{Variável Crítica} & \textbf{\% Ausente/Ignorado} & \textbf{Impacto Clínico} \\ \midrule
            Baciloscopia (\texttt{BACILOSCOP}) & \textbf{50,47\%} & Classificação Operacional \\
            Nervos Afetados (\texttt{NERVOSAFET}) & \textbf{47,29\%} & Predição de Incapacidade \\ \bottomrule
        \end{tabular}
        \caption{Qualidade dos dados brutos identificada na EDA.}
    \end{table}
\end{frame}

\begin{frame}{Dinâmica Temporal e a Subnotificação Pandêmica}
    \begin{columns}
        \column{0.6\textwidth}
        \begin{center}
            \includegraphics[width=\textwidth]{fig/comparativo_focado_ts.png}
        \end{center}
        
        \column{0.4\textwidth}
        \begin{itemize}
            \item \textbf{Gap Acumulado:} Identificamos uma perda de \textbf{33,56\%} no volume de detecções (2020-2022).
            \item \textbf{Recuperação (2023-2024):} Evidência clara de reabsorção de \textit{backlog} (supernotificação compensatória).
            \item \textbf{Conclusão:} O sistema está detectando o "passivo" gerado na pandemia.
        \end{itemize}
    \end{columns}
\end{frame}

\begin{frame}{Evolução Clínica: Matriz de Transição de Incapacidade}
    \begin{columns}
        \column{0.6\textwidth}
        \begin{center}
            \includegraphics[width=0.9\textwidth]{fig/matriz_transicao_gif.png}
        \end{center}
        
        \column{0.4\textwidth}
        \begin{itemize}
            \item \textbf{Inércia Clínica:} Observa-se que pacientes com Grau 2 no diagnóstico tendem a permanecer com deformidade.
            \item \textbf{Estabilidade (Diagonal):} A maioria dos pacientes fura o ciclo de reabilitação.
            \item \textbf{Impacto:} Reforça a necessidade do diagnóstico precoce.
        \end{itemize}
    \end{columns}
\end{frame}

% ============================================================================
% CAPÍTULO 3: METODOLOGIA
% ============================================================================
\section{Estratégia de Modelagem}

\begin{frame}{Divisão Temporal e Validação Walk-forward}
    Diferente da validação cruzada tradicional (K-Fold), adotamos o rigor da \textbf{Validação Walk-forward}:
    
    \begin{columns}
        \column{0.5\textwidth}
        \begin{center}
            \includegraphics[width=0.9\textwidth]{fig/comparativo_metodos_validacao.png}
        \end{center}
        
        \column{0.5\textwidth}
        \begin{itemize}
            \item \textbf{Motivação:} Evita o \textit{data leakage} (vazamento de dados do futuro para o passado).
            \item \textbf{Realismo:} Simula a expansão do conhecimento clínico conforme o tempo progride.
            \item \textbf{Baseline:} Dados pré-2020.
            \item \textbf{Teste:} Enfrentamento da ruptura pandêmica.
        \end{itemize}
    \end{columns}
\end{frame}

% ============================================================================
% CAPÍTULO 4: RESULTADOS
% ============================================================================
\section{Resultados e Discussão}

\begin{frame}{Performance do Modelo (LightGBM + SMOTE)}
    O LightGBM destacou-se pela eficiência com dados tabulares esparsos e desbalanceados.

    \begin{table}[ht]
    \centering
    \footnotesize
    \begin{tabular}{|l|c|c|c|c|}
    \hline
    \textbf{Métrica} & \textbf{Precision} & \textbf{Recall} & \textbf{F1-Score} & \textbf{AUC-ROC} \\ \hline
    Baixo Risco (G0) & 0.99 & 0.99 & 0.99 & \multirow{2}{*}{\textbf{0.93}} \\ \cline{1-4}
    \textbf{Alto Risco (G2D)} & \textbf{0.86} & \textbf{0.84} & \textbf{0.85} & \\ \hline
    \end{tabular}
    \caption{Métricas finais validadas no cenário pós-pandemia.}
    \end{table}
    
    \begin{itemize}
        \item \textbf{Alta Sensibilidade:} Fundamental para não ignorar pacientes que necessitam de intervenção imediata.
        \item \textbf{Robustez:} Manutenção da performance mesmo em dados com alta taxa de valores omissos.
    \end{itemize}
\end{frame}

\begin{frame}{Análise de Erros: Matriz de Confusão}
    \begin{columns}
        \column{0.5\textwidth}
        \textbf{Instrução para Figura 3:}
        \begin{itemize}
            \item \textit{Inserir a Matriz de Confusão do melhor modelo (LightGBM).}
            \item \textit{Destacar o quadrante dos Verdadeiros Positivos para G2D.}
        \end{itemize}
        % \includegraphics[width=\textwidth]{figuras/confusion_matrix_lightgbm.png}
        
        \column{0.5\textwidth}
        \begin{itemize}
            \item O modelo acertou 83\% dos casos graves.
            \item A "sujeira" nos dados (missing values) foi mitigada pela robustez do algoritmo de Gradient Boosting.
        \end{itemize}
    \end{columns}
\end{frame}

\begin{frame}{Justificativa da Rigidez Algorítmica}
    A escolha por UMAP e K=2 baseou-se em métricas de estabilidade e separabilidade:
    
    \begin{columns}
        \column{0.5\textwidth}
        \begin{center}
            \includegraphics[width=0.9\textwidth]{fig/comparativo_projecoes.png}
            \caption{UMAP preserva melhor a topologia clínica.}
        \end{center}
        
        \column{0.5\textwidth}
        \begin{center}
            \includegraphics[width=0.9\textwidth]{fig/otimizacao_silhouette.png}
            \caption{Silhouette Profile valida K=2 como ideal.}
        \end{center}
    \end{columns}
\end{frame}

\begin{frame}{Explicabilidade Global: Ranking de Importância (SHAP)}
    Além da visão individual (Waterfall), o modelo permite entender o comportamento sistêmico:
    
    \begin{columns}
        \column{0.6\textwidth}
        \begin{center}
            \includegraphics[width=\textwidth]{fig/shap_bar_global.png}
        \end{center}
        
        \column{0.4\textwidth}
        \begin{itemize}
            \item \textbf{Doses Recebidas:} Variável de maior impacto (proxie de adesão e tempo de tratamento).
            \item \textbf{Classificação Operacional:} MB vs PB define a probabilidade base de G2D.
            \item \textbf{Interpretabilidade:} O modelo não é uma caixa-preta; ele espelha diretrizes clínicas.
        \end{itemize}
    \end{columns}
\end{frame}

\begin{frame}{Perfís Clínicos via Clustering (UMAP + K-Means)}
    A clusterização revelou dois fenótipos populacionais distintos ($K=2$, Silhueta 0.69):
    
    \begin{columns}
        \column{0.55\textwidth}
        \begin{center}
            \includegraphics[width=0.9\textwidth]{fig/justificativa_alta_gravidade.png}
        \end{center}
        
        \column{0.45\textwidth}
        \begin{itemize}
            \item \textbf{Cluster 1 (Baixa Gravidade):} Diagnóstico precoce, maioria G0 (92,5\%).
            \item \textbf{Cluster 2 (Alta Gravidade):} Concentração massiva de \textbf{G2D (65,4\%)} e Multibacilares.
            \item \textbf{Utilidade:} Priorização de recursos para microrregiões no Cluster 2.
        \end{itemize}
    \end{columns}
\end{frame}

% ============================================================================
% CAPÍTULO 5: CONCLUSÃO
% ============================================================================
\section{Conclusão}

\begin{frame}{Conclusões e Impacto}
    \begin{alertblock}{O Passivo Oculto}
        A pandemia gerou um "apagão" seletivo. A análise de discrepância (Previsto vs. Observado) sugere que muitos casos G2D deixaram de ser notificados ou foram diagnosticados tardiamente.
    \end{alertblock}

    \begin{block}{Recomendação Tecnológica}
        \begin{itemize}
            \item O \textbf{LightGBM com SMOTE} provou ser viável para triagem em larga escala no SUS.
            \item Capacidade de lidar com dados incompletos (Baciloscopia ausente) é essencial para a realidade da Atenção Primária.
        \end{itemize}
    \end{block}
\end{frame}

% ============================================================================
% REFERÊNCIAS
% ============================================================================
\section{Referências}

\begin{frame}[allowframebreaks]{Referências Bibliográficas}
    \tiny
    \bibliographystyle{plainnat}
    \bibliography{referencias_dis} 
\end{frame}

\end{document}